%!TeX root = main.tex
%!encoding = UTF-8 Unicode

\section{Measurement-Driven Requirements for the Active Voltage Probe}
\label{sec:VoltageRequirements}

Prior to designing the dual-channel differential active voltage probe, its requirements are established. 
These requirements are specifically tailored for discrete SiC-MOSFETs, the primary focus of this work. 
Should the target devices change to SiC power modules or GaN devices, some modifications are necessary. 
The requirements are defined in accordance with the measurement limitations introduced in section \ref{sec:MeasurementLimitation}. 

\subsubsection*{Small Signal Response}
The small signal response of the probe is characterized by its bandwidth, which must exceed the maximum frequency of interest \mi{f}{interest,max} relevant to the measurements.
As discussed in section \ref{ch:methods}, this maximum frequency differs between low-pass and high-pass measurements.
For low-pass measurements, the bandwidth requirement corresponds to the maximum frequency of voltage components still above the noise floor.
For high-pass measurements, the maximum frequency of interest is set to \SI{500}{MHz}, which represents the point at which EMC-related issues transition from being caused by semiconductor switching events to those arising from digital signal sources, such as microcontroller and FPGA clock signals.

\subsubsection*{Large Signal Response}
The large signal response of the probe is characterized by its maximum input voltage, which must exceed the maximum voltage of interest \mi{V}{interest,max} relevant to the measurements.
Since the target of this work are \SI{1200}{V} SiC-MOSFETs, the maximum voltage of interest is set to \SI{1200}{V}.
Furthermore, the minimal rise time is determined by the minimal rise time possible in a given $LC$ resonant tank.
The stray inductance of the setup is around \SI{25}{nH}, and the output capacitance of the smallest investigated MOSFET, which is a \SI{30}{A} device is \SI{74}{pF}.
The rise time of an underdamped LC resonant tank is given by:

\begin{equation}
t_r = \frac{\pi}{\omega_n\sqrt{1-\zeta^2}} - \frac{\arcsin(\zeta)}{\omega_n\sqrt{1-\zeta^2}}
\end{equation}

where the natural frequency is $\omega_n = \frac{1}{\sqrt{LC}}$ and the damping ratio is $\zeta = \frac{R}{2}\sqrt{\frac{C}{L}}$ for a series RLC circuit.

For lightly damped systems ($\zeta \ll 1$), this simplifies to:

\begin{equation}
t_r \approx 1.8\cdot\sqrt{LC}
\end{equation}

This results in a minimal rise time of approximately \SI{2.5}{ns} for the given LC parameters.
However, this value represents a purely passive response that assumes an infinitely fast switch, a condition that is typically observed only during uncontrolled diode snap-off events.
In classical operation, the switching transient is predominantly determined by the gate-drive and is typically in the range of several tens of nanoseconds.


\subsubsection*{Parasitic Coupling}
As in any practical measurement setup, parasitic coupling can occur between the probe and the DUT.
For the proposed voltage probe, the dominant mechanism is capacitive coupling from the switching-node potential into the sensitive amplification stage.
Because this influence depends strongly on layout- and setup-specific conditions, it cannot be represented by a single precise quantitative requirement.
Therefore, the final measurement results must be validated by cross-checking against a trusted reference, typically a conventional passive voltage probe.

\subsubsection*{Transmission Line Effects}
In high-frequency measurements, transmission line effects can significantly impact the accuracy of the voltage probe.
This is explicitly a problem for voltage measurement, since the source impedance of the MOSFET is changing over the switching event and the load impedance of the probe is very high in range of ${M\Omega}$.
The transmission line impedance of the connection is normally in the range of \SI{50}{\Omega} to \SI{75}{\Omega}, which is significantly lower than the input impedance of the probe.
Hence, to avoid reflections and signal distortion, the connection between the probe and the DUT must be as short as possible, but maximally

\begin{align}
    \mi{l}{max} \leq \frac{\lambda_\mathrm{min}}{10} &= \frac{c_0}{10\cdot\mi{f}{max}\sqrt{\epsilon_r}} \\
                                                    &= \frac{c_0}{10\cdot \SI{500}{MHz}\cdot \sqrt{4.25}} \\
                                                    &= \SI{2.9}{cm}
\end{align}

This allowable distance is too short to be implemented in a practically useful setup with a cable connection.
Therefore, the probe must be designed for a direct connection to the drain and source terminals of the DUT.